\section{The remaining covariance and raw-summation interfaces}
\label{sec:closure-contracts}
We record the exact analytical implications needed in the low-frequency
and outer-frequency parts of the original density theorem. These
criteria distinguish proved finite-record facts from the two infinite
sums that still require estimates.

\begin{proposition}[A covariance closure criterion]
\label{prop:covariance-closure}
Put $g_R=G_R-\bar G_R$. Suppose that
\begin{equation}\label{eq:covariance-tail-contract}
 \lim_{m\to\infty}\sup_R\sum_{k>m}
       \left\|\int g_R\otimes(g_R\circ(F_R^*)^k)\dd\nu_R^*\right\|=0.
\end{equation}
Suppose also that zero asymptotic variance of $v\cdot g_R$ implies a
coboundary whose transfer function has a representative for which the
coboundary identity telescopes on the five real-rank periods of
Corollary~\ref{cor:real-period-rank}. Then the Green--Kubo matrices
\[
 \Sigma_R=\int g_R\otimes g_R\dd\nu_R^*
  +\sum_{k\ge1}\int\bigl(g_R\otimes(g_R\circ F^k)
                    +(g_R\circ F^k)\otimes g_R\bigr)\dd\nu_R^*
\]
are continuous and satisfy $cI_4\le\Sigma_R\le CI_4$ on $I$ for some
$c,C>0$.
\end{proposition}
\begin{proof}
Corollary~\ref{cor:all-moment-continuity} makes every fixed-lag matrix
continuous, including the zero-lag term. The uniformly summable tail in
\eqref{eq:covariance-tail-contract} therefore gives continuity of the
series. The usual expansion of the covariance of $\sum_{j<n}g_R\circ F^j$
is a finite identity; dividing it by $n$ and using absolute summability
identifies its limit with $\Sigma_R$. Thus $\Sigma_R$ is positive
semidefinite and its directional quadratic form is the asymptotic
variance.

If $v^T\Sigma_Rv=0$, the stipulated coboundary identity gives on each of
the five selected periods
\[
 v_1K_1+v_2K_2+v_3N+v_4T-h(v\cdot\bar G_R)=0.
\]
The augmented vector $(v_1,v_2,v_3,-v\cdot\bar G_R,v_4)$ annihilates
their five-dimensional record matrix, whose determinant is nonzero by
Corollary~\ref{cor:real-period-rank}. Hence $v=0$. Each matrix is
positive definite. Continuity and compactness of $I$ give uniform lower
and upper eigenvalue bounds. The periodic representative premise is
used explicitly here; a merely measurable identity is insufficient.
\end{proof}

\begin{lemma}[An absolute raw-residual summation criterion]
\label{lem:raw-residual-sum}
Let $\mathcal W$ be countable. For $w\in\mathcal W$ assign a lattice
label $k_w\in\Z^3$ and a function $q_w\in L^1(\R)$ whose second
distributional derivative is a finite signed measure. Suppose
\[
 A_0=\sum_w\|q_w\|_1<\infty,\qquad
 A_2=\sum_w\|D^2q_w\|_{\TV}<\infty.
\]
Then
$H(u,b)=\sum_w e^{iu\cdot k_w}\widehat q_w(b)$ belongs to
$L^1(\T^3\times\R)$ and, for $B\ge1$,
\begin{align}\label{eq:absolute-residual-sum}
 \int_{\T^3\times\R}|H|&\le2(2\pi)^3(A_0+A_2),\\
 \int_{\T^3\times\{|b|>B\}}|H|&\le2(2\pi)^3A_2/B.\nonumber
\end{align}
The same estimates hold uniformly for a parameter family whenever the
corresponding sums are uniformly bounded.
\end{lemma}
\begin{proof}
Distributional differentiation and the stated Fourier convention give
$-b^2\widehat q_w(b)=\widehat{D^2q_w}(b)$. Thus
\[
 |\widehat q_w(b)|\le\min\{\|q_w\|_1,
                    |b|^{-2}\|D^2q_w\|_{\TV}\}\quad(b\ne0).
\]
Integrate the first estimate on $|b|\le1$ and the second on $|b|>1$,
and then sum by Tonelli. The torus has volume $(2\pi)^3$.
The series itself converges uniformly as a transform of the absolutely
summable $L^1$ densities. Integrating only $|b|>B$ gives the second
bound. All lattice multiplicities are included in the two sums; no
bound on the number of words has been suppressed.
\end{proof}

To use this lemma, the residuals must be those of a decomposition of the
actual full measure, including singular itinerary boundaries. A jump left
in $q_w$ contributes a derivative of a point mass to $D^2q_w$ and does
not satisfy the hypothesis. The individual jump estimates of
Proposition~\ref{prop:general-edge} do not prove the sum $A_2<\infty$.
Nor does the scalar tail \eqref{eq:exponential-return} bound inverse
Jacobian factors or the variation of a density near a critical value.
If $A_2$ grows with $n$, its actual growth must be used in the frequency
splice rather than treated as a fixed constant.

There is nevertheless a useful exact truncation budget. For an insertion
$|w|\le M_0$, deleting the physical words with $N_{n,R}>L$ changes its
record measure, and its raw transform in uniform norm, by at most
\begin{equation}\label{eq:physical-count-truncation}
 M_0C_1e^{-a_1L/n}.
\end{equation}
Indeed total variation contracts under pushforward, and
\eqref{eq:block-exponential-moment} bounds the discarded mass.
On a frequency band of volume $V$ the integrated error is at most
$VM_0C_1e^{-a_1L/n}$. Thus the sufficient physical-count budget for a
prescribed integrated error $\varepsilon$ is
\[
 L\ge\frac{n}{a_1}\max\{0,\log(VM_0C_1/\varepsilon)\}.
\]
This controls a finite-band error, not a pointwise density error or an
infinite frequency band. It is compatible with retaining the original
raw datum and with the separate all-branch derivative contract above.
